\documentclass[../main.tex]{subfiles}

\begin{document}

%% ======================================================================
%%  Section 3 – HDF5 Core Library and File Format
%% ======================================================================
\section{HDF5 Core Library and File Format}
\label{sec:core-lib-format}

\subsection{HDF5 Safety Observations and Evidence Gaps}

Safety in this chapter means robustness and data integrity in the absence of a
malicious actor, as defined in Section~\ref{sec:assessment-dimensions}. This
revision examines four condition families and registers five formal findings:
two write-side findings, one filter-dependent archival finding, and two
reader/writer compatibility findings. A separate format-governance observation
remains in the evidence backlog because its adverse consequence and exposure
are not yet bounded.

The first three condition families are write-side and lifecycle risks and were the least developed part of
an earlier revision of this report; the discussion below states each as a
distinct trigger, because conflating them --- as an earlier draft did by
merging storage exhaustion and abrupt termination into one observation ---
obscures that they need different containment and different structural
treatment.

\begin{description}[leftmargin=1.5em,style=nextline]
    \item[Storage exhaustion during write or close.]
    Storage fills during a write, flush, or close, and the operation reports
    failure through the ordinary HDF5 error path. The library offers no
    structured VFD-level notification at the failure point, defined recovery
    procedure, or supported way to mark the file suspect and abandon it safely.
    The controlling evidence is a maintainer statement that there is no HDF5
    recovery story for the condition and no VFD callback for it, together with
    four candidate treatments and a documented cache workaround used in its
    place~\cite{forum-enospc-11714}. Registered as
    \code{CORE-SAFE-001}, \textbf{Moderate--High}. An earlier revision recorded
    this as NE for want of a bounded recovery path; the recovery path is in
    absent at the cited architectural boundary. That establishes a control gap;
    it does not establish how often storage exhaustion occurs or what durable
    state follows it.

    \item[Corruption from abrupt termination.]
    A writer is killed, crashes, or loses power with metadata updates in
    flight, and in-memory and durable state diverge with no journal or
    write-ahead log to reconcile them. Community reports recur across sixteen
    years, and The HDF Group's own 2024 discussion states that a crash ``can
    lead to data loss or corruption of the file'' and describes restarting
    crash-proofing work
    \cite{forum-recover-1146,forum-crash-state-1598,forum-crash-prevent-2462,
    forum-flush-crash-3481,forum-flush-corrupt-4958,forum-recover-5441,
    forum-crash-corruption-7581,forum-crashproofing-12003}. Registered as
    \code{CORE-SAFE-002}, \textbf{Moderate--High}. This condition had no record
    of any kind in an earlier revision.

    \item[Extension-dependent loss of interpretability.]
    An archived dataset requires a non-core filter that is later unavailable.
    Registered as \code{LONG-SAF-001}, \textbf{Moderate} at
    a stated ten-year-or-more horizon, and developed in
    Section~\ref{sec:long-haul}. Because the horizon is material, decision
    rule 7 forbids ordering this rating directly against ratings that state no
    horizon. VFD and VOL archival dependencies remain a reasoned extension of
    the inquiry, not part of this rated scenario.

    \item[Reader/writer disagreement.]
    The two 2026 chunk-layout cases in Section~\ref{sec:indie}, registered as
    \code{IND-SAF-001} and \code{IND-SAF-002}. They share an unreadability
    outcome with \code{LONG-SAF-001}, but not its mechanism: these are
    specification and metadata-conformance failures, whereas the archival
    finding concerns an unavailable runtime dependency.
\end{description}

\noindent\textbf{The shape of the gap this closes.}
It is worth naming why safety was the thinnest dimension. The rest of this
report is organised around \emph{read-side hostile input} --- the rubric's
principal threat scenario, the CVE analysis, the invariant registry, the parser
boundary, the fuzzing programme, and the security profiles all answer the
question ``a malicious file has arrived.'' The three findings above answer a
different question, ``my file did not survive,'' and no structural
recommendation in the security track addresses it. That asymmetry, rather than
any judgment that the write path is sound, is why these conditions were
previously unregistered or unrated.

\noindent\textbf{Coverage limitation, stated against an existing hazard
analysis.}
Safety is the least developed dimension in this revision. The
extent of the gap can be stated precisely rather than vaguely, because the SSP
SIG has already published a sixteen-entry safety hazard registry built on a
control-system hazard model with unsafe control actions, derived safety
constraints, and per-entry
ratings~\cite{hdf5-ssp-haz,hdf5-safety-hazard-model}.

Table~\ref{tab:haz-crosswalk} maps that registry to this report.
\textbf{Of sixteen hazards, this revision registers a corresponding formal
finding for four, an unrated evidence record for two, and nothing at all for
ten.} An earlier revision covered one. That is the honest measure of this chapter's safety coverage. The
registry entries are not reproduced or re-rated here --- their method is
incompatible with Section~\ref{sec:risk-rubric}, as
Section~\ref{sec:rubric-reconciliation} explains --- and where this report is
silent, the registry entry stands as the SSP SIG's own assessment. Silence
here is not disagreement and is emphatically not a Low-risk conclusion.

\begingroup
\small
\setlength{\tabcolsep}{3pt}
\renewcommand{\arraystretch}{1.15}
\setlength{\LTleft}{-0.9in}
\setlength{\LTright}{-0.9in}
\begin{longtable}{@{}>{\raggedright\arraybackslash}p{0.09\widetablewidth}>{\raggedright\arraybackslash}p{0.31\widetablewidth}>{\raggedright\arraybackslash}p{0.11\widetablewidth}>{\raggedright\arraybackslash}p{0.45\widetablewidth}@{}}
\caption{SSP SIG safety hazard registry mapped to this report's coverage}
\label{tab:haz-crosswalk}\\
\hline
\textbf{ID} & \textbf{Hazard} & \textbf{Their rating} &
\textbf{Coverage in this report} \\
\hline
\endfirsthead

\caption[]{SSP SIG safety hazard registry mapped to this report's coverage,
continued}\\
\hline
\textbf{ID} & \textbf{Hazard} & \textbf{Their rating} &
\textbf{Coverage in this report} \\
\hline
\endhead

\hline
\endfoot

\hline
\endlastfoot

\code{HAZ-001} & Crash during metadata update & 20, Very High &
\textbf{Covered.} \code{CORE-SAFE-002}, Moderate--High, Near-term containment
with structural treatment under \code{CORE-SAFE-S2}. Registered in this
revision; absent from the previous one. \\

\code{HAZ-002} & Concurrent access without safe coordination & 15, High &
\textbf{Not assessed.} No concurrency scenario was developed. \\

\code{HAZ-003} & Extension boundary violates core invariants & 15, High &
\textbf{Partially covered as an unrated record.} The privacy aspect is
\code{EXT-PRIV-001}; the \emph{safety} aspect --- an extension breaking an
invariant the core relies on --- has no counterpart here. \\

\code{HAZ-004} & Metadata and artifacts leak sensitive information & 12, High &
\textbf{Covered.} \code{CORE-PRIV-001}, Moderate--High, Near-term. This is the
one hazard with a corresponding rated finding. \\

\code{HAZ-005} & Out-of-space during write or close & 15, High &
\textbf{Covered.} \code{CORE-SAFE-001}, Moderate--High, Near-term containment
with structural treatment under \code{CORE-SAFE-S1}. The registry's accident
chain and safety constraints SC-1 to SC-3 are adopted as the \code{RM-2D}
hypotheses. \\

\code{HAZ-006} & False sense of durability from flush & 12, High &
\textbf{Not separately assessed}, but within the evidence base for
\code{CORE-SAFE-002}: the community record includes reports that flushing
itself can leave a file corrupted. A distinct durability-signaling finding is
still warranted. \\

\code{HAZ-007} & SWMR deployment misfit & 12, High &
\textbf{Not assessed.} \\

\code{HAZ-008} & File-locking pitfalls or disabling via environment variable &
12, High &
\textbf{Not assessed.} File locking appears only as one item in the
\code{SUP-SAFE-001} mismatch list, which is itself unrated. \\

\code{HAZ-009} & Thread-safe build uses a global lock & 12, High &
\textbf{Not assessed.} \\

\code{HAZ-010} & Chunk size too small & 12, High &
\textbf{Not assessed as a safety hazard.} Pathological chunk shapes appear
here only as an adversarial resource-consumption mechanism under
\code{CORE-SEC-002}; the non-adversarial operational hazard is absent. \\

\code{HAZ-011} & Deleting data does not shrink files by default & 12, High &
\textbf{Not assessed.} Note the privacy adjacency: the registry separately
rates data remanence after deletion as \code{EXP-003}. \\

\code{HAZ-012} & Recovery tooling limits & 12, High &
\textbf{Not separately assessed.} Recovery capability is a controlling
assumption of the \code{CORE-SAFE-002} impact range --- whether the file
reopens after \code{h5clear} is what separates I2 from I3 --- so the limits of
that tooling bear directly on the rating. The evidence-gated repair planner in
H5Lens is relevant prior art (Section~\ref{sec:core-prior-art}). \\

\code{HAZ-013} & High-level API thread-safety gaps & 9, Moderate &
\textbf{Not assessed.} \\

\code{HAZ-014} & Portability hazard from missing filter plugins & 9, Moderate &
\textbf{Covered.} \code{LONG-SAF-001} rates the missing-filter archival
scenario at a stated decadal horizon. VFD and VOL archival dependencies are
identified as unmeasured extensions of the inquiry, not as part of the rating. \\

\code{HAZ-015} & Misuse of metadata-cache flush controls & 8, Moderate &
\textbf{Not assessed.} \\

\code{HAZ-016} & Chunk cache mis-sizing & 6, Moderate &
\textbf{Not assessed.} \\

\end{longtable}
\endgroup

\noindent\textbf{Consequence for the reader and for the roadmap.}
Three things follow. First, a reader must not treat this report's small number
of safety findings as evidence that HDF5's safety surface is small; it is
evidence that this engagement examined it least. Second, the natural next
revision of this report should adopt the hazard registry as its safety
enumeration and work through it scenario by scenario, rather than re-deriving a
list. Third, \code{RM-2D} is currently scoped to storage exhaustion alone; the
registry shows at least \code{HAZ-001}, \code{HAZ-006}, \code{HAZ-012}, and
\code{HAZ-015} share its evidence apparatus --- fault injection at write,
flush, and close boundaries, followed by durable-state and recovery inspection
--- and it has accordingly been widened to cover the whole cluster, including
the abrupt-termination trigger behind \code{CORE-SAFE-002}. Ten hazards remain
with no record here; that number, not the four now covered, is the honest
measure of what a safety-focused revision would still have to do.

\subsection{Malformed Input as a Security Boundary}

Three mechanisms recur across the security evidence for the core library and
official tools: arithmetic on file-derived quantities that is not checked for
overflow or zero denominators; file-format parsing that constructs internal
objects before validating the bytes that describe them; and executable
extension code reached through plugin loading, which is treated in
Section~\ref{sec:extensions}. The first two are examined here.

A CVE history summary is included in Appendix~\ref{app:cve-history-summary}.
The list mixes real HDF5-library issues, removed-tool issues,
duplicate/rejected/not-HDF5 entries, and entries with no commit URL.
A recurring class represented in the selected history is failure to validate
file-derived structure sizes, offsets, counts, message lengths, continuation
chunks, datatype metadata, or filter parameters before using them in low-level
C memory operations. The 2025 entries include heap overflows, use-after-free,
null dereference, memory leak, and resource-consumption defects across object
headers, file-space metadata, type conversion, filters, and metadata cache
code. This qualitative clustering does not estimate ecosystem frequency.

The fix commentary confirms that several CVEs were symptoms of corrupted metadata
not being rejected early. (Examples: the file-space page-size fix added a sanity
check immediately after reading \code{page\_size} from the file; the
object-header-continuation fix rejects zero continuation chunk sizes;
the self-referential continuation-message fix checks expected versus actual
deserialized object-header chunks; and the cache-image fix added buffer-size
arguments, overflow checks, input validation, and cleanup.)

An August 2026 source-review observation exposes a related but distinct control
gap. HDF5 Pickles issue~\#87 records candidate deserializer invariants whose
truth depends on in-file bytes but whose checks are masked inside
\code{assert(...)} expressions; those expressions provide no runtime validation
when assertions are disabled with \code{-DNDEBUG}. The open catalog task names
version-2 B-trees, dataset chunk records, extensible arrays, fractal heaps,
free-space managers, shared object-header messages, and metadata-cache images,
with a cursory estimate of 50--100 instances
~\cite{hdf5-pickles-issue-87}. A fixed August 25 develop snapshot corroborates
the mechanism: its build configuration notes that CMake adds \code{NDEBUG} for
Release builds, while representative source sites assert a decoded filtered-
chunk address and byte count before returning success, assert extensible-array
parameters before geometry and allocation calculations, and assert a
reverse-filter output size before copying the expected fractal-heap block size
~\cite{hdf5-assert-build-config-2026,hdf5-assert-chunk-record-2026,
hdf5-assert-extensible-array-2026,hdf5-assert-fractal-heap-2026}. This static
sample confirms a source-level assurance gap, but the 50--100 estimate is not a
verified site count, and the available evidence does not establish that every
candidate is a sole guard, release-reachable in a supported version, or capable
of a particular adverse outcome. This report therefore tracks the condition as
evidence record \code{CORE-SEC-003}, not as a separately rated formal finding.

One member of that population has been carried further, and it is worth stating
precisely because it illustrates both the mechanism and the remaining evidence
boundary. A version-2 B-tree header declaring a zero record size reaches an
integer division while the leaf record capacity is derived; the sole guard is
an assertion in \code{H5B2hdr.c}, absent under \code{-DNDEBUG}. The cited case
record reports \code{SIGFPE} from a locally built \code{h5dump -H} and the
public \code{H5Fopen}/\code{H5Gopen2}/\code{H5Gget\_info} path using a
release-configured HDF5 2.3.0 development source tree
~\cite{h5policy-case-v2btree-recordsize}. It expressly records no 32-bit
measurement and does not establish an affected supported release or version
range. The case therefore corroborates an adverse \code{-DNDEBUG} path but
does not satisfy this report's supported-version condition for promotion into
a formal finding. It remains evidence under \code{CORE-SEC-003}. A second
candidate, the cache-image LRU rank relation, is reported without a measured
release-build consequence~\cite{h5policy-case-mdci-lru}. This bounded
reproduction is a reason to complete the supported-release catalog, not a
substitute for it.
Regardless of the final count, a predicate whose validity depends on file bytes
must have an always-active runtime check; an assertion may duplicate that check
for diagnostics but must not enforce it alone.

The observed root-cause categories are more compact than a CVE-by-CVE list and
are summarized in Table~\ref{tab:hdf5-cve-root-cause-categories}. They were
derived independently in this assessment, and they are consistent with the SSP
SIG security threat model, which decomposes the same evidence into thirteen
mechanism entries~\cite{hdf5-security-threat-model,hdf5-ssp-atk}; that
decomposition is the finer-grained one and is the better basis for the
invariant tranches proposed in \code{CORE-S1}.

Table~\ref{tab:root-cause-cwe} maps each category to the standard weakness
vocabulary, using the SSP SIG's top-CWE list for HDF5
contexts~\cite{hdf5-top-cwe}. The mapping is worth stating explicitly because
it makes the report's central security claim checkable against an external
taxonomy rather than only against its own categories: nearly every entry
resolves to \code{CWE-20} and its children as the \emph{cause}, with the
memory-safety CWEs appearing as \emph{consequences}. Per-CVE tags are in
Appendix~\ref{app:cve-history-summary}.

\begingroup
\small
\setlength{\tabcolsep}{3pt}
\renewcommand{\arraystretch}{1.15}

\setlength{\LTleft}{-0.7in}
\setlength{\LTright}{-0.7in}

\begin{longtable}{@{}p{0.28\widetablewidth}p{0.68\widetablewidth}@{}}
\caption{Observed root-cause categories in the selected HDF5 CVE history}
\label{tab:hdf5-cve-root-cause-categories}\\
\hline
\textbf{Root-cause category} & \textbf{What it means for HDF5} \\
\hline
\endfirsthead

\caption[]{Observed root-cause categories in the selected HDF5 CVE history, continued}\\
\hline
\textbf{Root-cause category} & \textbf{What it means for HDF5} \\
\hline
\endhead

\hline
\endfoot

\hline
\endlastfoot

unchecked file-derived sizes / counts / offsets &
A recurring evidenced class. HDF5 metadata controls object layout, heap blocks, chunks,
dataspaces, filters, links, and type descriptors. When those fields are malformed,
C code later trusts them. \\
inconsistent internal bookkeeping &
Attribute tables, cache images, continuation chunks, and free lists often needed
separate ``allocated capacity'' versus ``actual count'' tracking. \\
decode/encode asymmetry &
Several issues happen because decode did not reject bad messages, then encode/flush
later assumed valid native structures. \\
arithmetic precondition failures &
Divide-by-zero and multiplication-overflow checks were missing in chunking, layout,
b-tree, filter, and dataspace paths. \\
unsafe cleanup/lifetime paths &
Use-after-free, double-free, leaks, and null dereferences show that malformed input
often breaks unwinding logic as much as the primary decode logic. \\
unbounded traversal, recursion, loop-progress, and resource-limit failures & 
File-controlled graphs, dimensions, and metadata can drive recursion, repeated
work, or allocation without a structural progress rule or explicit budget.\\
legacy tool attack surface &
\code{gif2h5}, HDF to GIF conversion, \code{h5dump}, and \code{h5repack} appear
repeatedly. Removal of \code{gif2h5} was a structural mitigation, not just a patch. \\
triage quality, duplicate/variant clustering, reachability, and fix-completeness metadata &
Several entries have \code{commit_url: null}, ``not filed against HDF5,'' ``cannot
reproduce,'' ``not applicable,'' or tentative confidence. Those should not be counted
as reviewed HDF5 fixes without qualification. \\
\end{longtable}
\endgroup

\begingroup
\small
\setlength{\tabcolsep}{3pt}
\renewcommand{\arraystretch}{1.15}
\setlength{\LTleft}{-0.9in}
\setlength{\LTright}{-0.9in}
\begin{longtable}{@{}>{\raggedright\arraybackslash}p{0.28\widetablewidth}>{\raggedright\arraybackslash}p{0.24\widetablewidth}>{\raggedright\arraybackslash}p{0.44\widetablewidth}@{}}
\caption{Observed root-cause categories mapped to standard weaknesses}
\label{tab:root-cause-cwe}\\
\hline
\textbf{Root-cause category} & \textbf{Cause (CWE)} &
\textbf{Consequence (CWE) and note} \\
\hline
\endfirsthead
\caption[]{Observed root-cause categories mapped to standard weaknesses,
continued}\\
\hline
\textbf{Root-cause category} & \textbf{Cause (CWE)} &
\textbf{Consequence (CWE) and note} \\
\hline
\endhead
\hline
\endfoot
\hline
\endlastfoot

unchecked file-derived sizes, counts, offsets &
\code{CWE-20}, \code{CWE-1284}, \code{CWE-1285} &
\code{CWE-787}, \code{CWE-125}. The single largest group. A quantity or an
offset taken from the file is used before it is validated. \\

inconsistent internal bookkeeping &
\code{CWE-20}, \code{CWE-1284} &
\code{CWE-787}, \code{CWE-908}. Allocated capacity and actual count are
conflated, so a validated bound does not describe the buffer actually held. \\

decode/encode asymmetry &
\code{CWE-20}, \code{CWE-1286} &
\code{CWE-787}. Decode accepts a structure that encode or flush later assumes
is well-formed; the syntactic check that was never made surfaces on the write
path. \\

arithmetic precondition failures &
\code{CWE-20}, \code{CWE-1284} &
\code{CWE-369}, \code{CWE-190}. Divisors and multiplications derived from file
bytes without a zero or overflow precondition. \\

unsafe cleanup and lifetime paths &
\code{CWE-20} &
\code{CWE-416}, \code{CWE-415}, \code{CWE-476}, \code{CWE-401}. Distinctive in
that the unwinding path, not the decode path, is what malformed input breaks. \\

unbounded traversal, recursion, and resource limits &
\code{CWE-20}, \code{CWE-1284} &
\code{CWE-400}, \code{CWE-674}. File-controlled graphs and counts drive work
with no structural progress rule or explicit budget. \\

legacy tool attack surface &
\code{CWE-20} &
\code{CWE-787}, \code{CWE-125}. Not a distinct weakness class; a distinct
\emph{exposure} class, since these are the programs users run first on
untrusted files. \\

type-confusion on message and datatype flags &
\code{CWE-20}, \code{CWE-1287} &
\code{CWE-787}. A message or datatype is cast to a structure its declared type
does not support --- validation of a specified \emph{type}, distinct from
validating a quantity or an index. \\

triage and fix-completeness metadata &
\emph{n/a} &
\emph{n/a}. A records-quality category, not a weakness. It affects what can be
concluded from the history, not what the library does. \\

\end{longtable}
\endgroup

Read together, the two tables state the report's core security claim in
standard terms: the memory-safety weaknesses that dominate the CVE history are
consequences, and the cause they share is \code{CWE-20}. That is why
\code{CORE-S1} through \code{CORE-S3} target validation and checked arithmetic
rather than the individual crash classes, and why closing a \code{CWE-787}
instance does not close the class that produced it.

\subsection{HDF5 Privacy: Assets and Exposure}

\subsubsection{Introduction} 

HDF5 design features create predictable privacy-exposure mechanisms, but
realized harm, exposure, and any lifecycle trend are deployment-specific. The
file format and library are one source of exposure surfaces; scientific
workflows, software ecosystems, interoperability mechanisms, and changing data
contexts can add others.

Before describing the assets that require protection in HDF5 and the mechanisms through
which privacy exposure may occur, it is important to emphasize several fundamental
characteristics of HDF5 and the audit approach adopted in this document.

\begin{itemize}
    \item \emph{HDF5 privacy primarily concerns accidental or unintended exposure of sensitive information
          rather than malicious attacks}. Malicious attacks and unauthorized access are primarily
          addressed through cybersecurity mechanisms and are outside the scope of this section.
    \item \emph{HDF5, by design, does not provide privacy protections}. When privacy is a concern, users and organizations
           are responsible for protecting data throughout the entire data lifecycle. In this document,
           the term data refers not only to individual HDF5 objects within a single file, but also to
           collections of HDF5 files and external data accessible through the HDF5 library,
           connectors, and related software components.
    \item \emph{HDF5 neither implements nor enforces privacy standards}. Furthermore, even when privacy
          standards and technical safeguards are in place, sensitive information may still be exposed
          through legitimate data processing, contextual interpretation, pattern discovery, inference,
          aggregation, and other analytical activities.
    \item \emph{HDF5 does not manage authentication, authorization, or filesystem protection}. However, privacy
          risks extend beyond traditional security concerns. Sensitive information may be exposed during
          legitimate data use by authorized users, even within highly secure computing environments.
    \item \emph{Additional lifecycle stages can add privacy-exposure surfaces}.
          Sensitive information stored in HDF5 may be exposed during data creation,
          processing, storage, sharing, transfer, reuse, or archival; classification,
          minimization, access restrictions, and lifecycle controls can increase or
          reduce the realized deployment risk.
    \item \emph{HDF5 extensibility mechanisms and associated software ecosystems may alter the privacy characteristics of data}.
          Features such as filters, Virtual Object Layer (VOL) connectors, Virtual File Drivers (VFDs), and
          external applications, including conversion and analysis tools, may introduce new privacy
          exposure surfaces or modify existing ones.
\end{itemize}

The remainder of this section focuses on the assets that require protection in HDF5, the mechanisms
through which sensitive information may be exposed, and potential approaches for mitigating privacy risks.

\subsubsection{HDF5 Assets and their Exposure}

This subsection identifies five classes of asset that may require protection
in an HDF5 deployment: the scientific data themselves, user-defined metadata,
structural metadata, information derived or reconstructed from combinations of
the former, and artifacts produced by the surrounding tools and workflows.
Each is treated in turn below. The classes are cumulative rather than
alternative --- a single sharing boundary can disclose through all five --- and
the enumeration in \code{CORE-PRIV-S1} expects a disposition for every class,
including a positive statement where a class was examined and found empty.

\paragraph{Data Assets}\mbox{}\\
The primary assets requiring protection are the scientific data themselves. In HDF5, these 
include data stored in HDF5 objects, such as files, groups, datasets, datatypes, and dataspaces, 
which represent experimental observations, measurements, simulations, images, 
sensor outputs, and computational results.

Scientific data may be exposed through direct access using HDF5-based software, including 
command-line utilities, graphical viewers, and programming interfaces such as Python. 
However, privacy-sensitive information may also be inferred through examination of the low-level 
storage organization and auxiliary metadata embedded within the file to support portability, 
interoperability, and efficient data access. For example, string attributes visible through 
standard operating system utilities (e.g., the Unix strings command) may disclose dataset 
names, experimental identifiers, or other descriptive information. Similarly, inspection 
of dataset headers may reveal storage layout, chunk organization, compression methods, 
or data acquisition patterns (e.g., frame-based acquisition stored as chunked datasets). 
Examination of embedded strings may also disclose the use of external storage or references 
to datasets located in other HDF5 files.

It is important to distinguish between protecting access to data values and protecting 
information about data representation. Even when raw data are protected through mechanisms 
such as compression or encryption, information describing their storage organization, layout, 
or access methods may remain accessible and reveal sensitive contextual information about 
the underlying data or the processes that generated them.


\paragraph{Metadata Assets}\mbox{}\\
User-defined metadata often constitute one of the most privacy-sensitive asset classes in HDF5. 
The format was intentionally designed as a self-describing data model that allows extensive 
contextual information to be stored alongside scientific data. While HDF5 attributes are the most 
recognizable form of user-defined metadata, they represent only one category of descriptive 
information maintained within an HDF5 file.

User-defined metadata include all user-created HDF5 objects and object properties that organize, 
identify, or describe stored data. These include groups, datasets, their associated datatypes 
and dataspaces, committed datatypes, and attributes attached to HDF5 objects. 
Each HDF5 attribute is itself a structured metadata object with a name, datatype, and dataspace. 
Additional metadata assets include file and object names (link names), committed datatype names, 
compound datatype member names, enumeration labels, object comments, and information 
stored in the user block. Collectively, these metadata describe the logical organization, 
semantics, and relationships of the stored data and frequently contain experiment-specific, 
operational, or organizational information.

Like scientific data, metadata assets may be exposed not only through HDF5 applications but 
also through direct examination of the binary file format. Inspection of file structures or 
extraction of embedded strings without using HDF5 software may reveal object names, 
descriptive attributes, storage relationships, user block contents, and other contextual 
information that could disclose sensitive information about the data, their origin, 
or the organization that produced them.

\paragraph{Structural Metadata Assets}\mbox{}\\
Privacy-sensitive information may also be inferred from the structural characteristics of HDF5 files. 
Structural metadata assets include user-defined organizational structures, such as 
the hierarchical arrangement of objects, dataset storage layouts (e.g., contiguous, chunked, 
compact, or external storage), filter pipelines, and relationships established through links, 
object references, region references, or virtual datasets (VDS). In addition, the HDF5 library 
maintains internal structural metadata that describe and manage these user-defined structures, 
including chunk indexing schemes, B-tree organizations, and other implementation-specific data structures.

Even when raw data are inaccessible, for example, because they are encrypted, compressed, 
stored externally, or protected by file-system access controls, structural metadata may 
disclose information about the purpose of the data, experimental design, workflow organization, 
relationships among datasets, or assumptions made by the generating application. 
Naming hierarchies, grouping strategies, storage layouts, and reference relationships may reveal 
semantic information about scientific activities or operational processes. Likewise, 
implementation-specific structures, such as chunk indexing formats or other internal metadata, 
may disclose characteristics of the HDF5 library version or software implementation 
used to create or process the file.

Like scientific data and user-defined metadata, structural metadata may be exposed through 
HDF5 applications or by direct examination of the binary file format without using HDF5 software.


\paragraph{Derived and Reconstructed Information}\mbox{}\\
Privacy-sensitive information in HDF5 may also arise through inference or reconstruction 
from combinations of available data and metadata. Derived assets include relationships 
inferred from multiple HDF5 datasets, external links, object references, 
region references, virtual dataset (VDS) mappings, or information obtained by 
combining HDF5 data with external sources. 
These relationships may reveal information that is not explicitly represented within any single dataset or file.

For example, correlating quality-control datasets with experimental observations, combining 
temporal and spatial measurements, or reconstructing relationships across linked HDF5 files 
may disclose sensitive characteristics that are not apparent when individual datasets 
are examined in isolation. Consequently, privacy-sensitive assets extend beyond explicitly 
stored data and metadata to include information that can be inferred through aggregation, 
correlation, reconstruction, or other forms of analysis.

\paragraph{Workflow and Applications Artifacts}\mbox{}\\
Privacy-sensitive assets in HDF5 extend beyond data and metadata stored within HDF5 files 
to include artifacts created by applications, processing workflows, and the surrounding 
software ecosystem. These artifacts include command outputs (e.g., h5dump output), 
diagnostic logs, temporary files, metadata snapshots, cached intermediate datasets, 
conversion artifacts, and outputs generated by plugins or connectors 
(e.g., binary files produced by an HDF5 VOL connector for another scientific data format).

Examples include VFD SWMR metadata snapshots, cached intermediate results, diagnostic logs, 
and metadata generated during translation between non-self-describing data formats and HDF5. 
Such artifacts may expose scientific data, metadata, storage organization, workflow decisions, 
or implementation details beyond what is contained in the original HDF5 file. 
Because they are often created automatically, these artifacts may persist beyond their 
intended lifetime and remain unnoticed by users focused primarily on the scientific datasets themselves.

Scientific data processing workflows routinely enrich, restructure, annotate, and transform 
datasets to improve interoperability, reproducibility, usability, or analytical value. 
However, these processes may also introduce or preserve contextual information that was not intended for disclosure. 
Automatically generated metadata, intermediate processing products, and derived datasets may 
reveal experimental assumptions, workflow decisions, data provenance, software 
implementation details, or scientific semantics that were absent from the original data. 
Consequently, privacy exposure may arise not from malicious activity or unauthorized access, 
but as an unintended consequence of routine scientific data processing.


\subsection{HDF5 Tools}

Not strictly part of the core library, but the tools are an important part of the ecosystem and have their own set of issues.

\subsubsection{Privacy exposure through HDF5 command-line tools}

HDF5 command-line tools provide convenient mechanisms for managing HDF5 files. 
They can be used to display or extract raw data, and user-defined and structural 
HDF5 metadata (e.g., \code{h5dump, h5ls, h5debug}), edit raw data and attributes with \code{h5edit}, 
extract or aggregate data using \code{h5copy}, change characteristics of HDF5 objects, 
for example, remove or add data filters and change storage properties using \code{h5repack}.
 HDF5 command-line tools are routinely employed throughout the HDF5 data lifecycle. 
While the command-line tools are essential for working with HDF5 data, they are ultimate 
sources of unintended privacy exposure.

The HDF5 tools were designed to ease access to HDF5 data, to reveal information about 
the contents and organization of HDF5 files, and to manage that information while 
avoiding steep learning curve of HDF5 programming. I.e., HDF5 command-line tools were 
designed to expose HDF5 assets. Depending on the tool and the selected options, they may 
display object names, attributes, datatype definitions, storage layouts, chunking parameters, 
filter pipelines, external storage locations, object references, virtual dataset mappings, 
and portions or all of the stored data to name a few. Even when users intend only to 
verify file integrity or inspect metadata, command output may disclose information that 
is considered sensitive in a particular context.

Privacy exposure is not limited to information displayed on the terminal by \code{h5dump} or \code{h5ls}. 
Command output is frequently redirected to files, incorporated into automated workflow logs, 
attached to bug reports, archived as diagnostic information, or shared. The output can be used by 
other tools (e.g., \code{h5import}) or applications to generate new data. Consequently, 
metadata and scientific data extracted by HDF5 tools may persist outside the original HDF5 
file and become accessible to individuals who were never intended to access the 
underlying information.

Several command-line tools may also create new HDF5 files or transform existing ones. \code{h5repack}
 and \code{h5copy} may generate additional copies of data and/or aggregate data, modify 
storage layouts, remove or add data filtering; \code{h5edit} 
may modify and introduce new metadata, change raw data values and preserve information 
that was not expected to be retained. I.e., the tools can alter the privacy 
characteristics of the resulting files, particularly when they are incorporated into 
larger workflows, backups, or public data repositories.

Because HDF5 command-line tools are widely used in automated scripts, continuous 
integration systems, and scientific workflows, their privacy implications should be 
evaluated alongside those of the HDF5 library and HDF5-based applications. 
The privacy mitigation strategies discussed in the next section apply not only to the 
HDF5 library itself but also, and often to an even greater extent, to HDF5 command-line tools. 
This is because these tools are designed to expose HDF5 assets and to generate 
persistent outputs (e.g., new binary and text files). 

The table below summarizes the assets exposed by the listed command-line tools.

\begingroup
\small
\setlength{\tabcolsep}{3pt}
\renewcommand{\arraystretch}{1.15}

\setlength{\LTleft}{-0.9in}
\setlength{\LTright}{-0.9in}

\begin{longtable}{@{}p{0.12\widetablewidth}p{0.22\widetablewidth}p{0.18\widetablewidth}p{0.48\widetablewidth}@{}}
\caption{Privacy exposure through HDF5 command-line tools}
\label{tab:hdf5-tools-privacy}\\
\hline
\textbf{Tool} &
\textbf{Exposed assets} &
\textbf{Workflow artifacts} &
\textbf{Privacy considerations} \\
\hline
\endfirsthead

\caption[]{Privacy exposure through HDF5 command-line tools, continued}\\
\hline
\textbf{Tool} &
\textbf{Exposed assets} &
\textbf{Workflow artifacts} &
\textbf{Privacy considerations} \\
\hline
\endhead

\hline
\endfoot

\texttt{h5dump} &
Raw data; user-defined and structural metadata &
Text dumps, redirected output, logs &
Displays datasets, attributes, datatypes, storage layout, filter information, object references, and VDS mappings. May expose data stored in other HDF5 files through external links and Virtual Datasets (VDS). Output is frequently redirected to files or logs, creating persistent copies of sensitive information.\\
\hline

\texttt{h5ls} &
User-defined and structural metadata (optionally raw data) &
Terminal output, logs &
Reveals object hierarchy, object names, datatypes, storage properties, links, and object relationships. With the \texttt{-d} option, also displays dataset values.\\
\hline

\texttt{h5repack} &
Raw data; user-defined and structural metadata &
New HDF5 file &
Creates additional HDF5 copies while potentially modifying storage layout, filters, and metadata. The \texttt{--merge} option can aggregate datasets reachable through external links into a single file. May access non-HDF5 data through configured VFDs and VOL connectors.\\
\hline

\texttt{h5diff} &
Raw data; user-defined metadata &
Comparison reports &
Reads and compares datasets and attributes. Reports may disclose sensitive values or structural differences.\\
\hline

\texttt{h5copy} &
Raw data; user-defined and structural metadata &
Additional HDF5 files &
Copies objects within or between HDF5 files, increasing the number of persistent copies of sensitive information.\\
\hline

\texttt{h5stat} &
Structural metadata &
Statistics reports, logs &
Reports object counts, storage statistics, chunk information, free space, and other implementation characteristics that may reveal dataset complexity or workflow properties.\\
\hline

\texttt{h5debug} &
Structural metadata &
Debug output &
Displays low-level internal file structures, addresses, B-trees, heaps, chunk indexes, and other implementation metadata intended for debugging.\\
\hline

\texttt{h5watch} &
Raw data; user-defined metadata &
Monitoring output, logs &
Monitors SWMR datasets and reports newly written data, timestamps, and update activity in real time.\\
\hline

\texttt{h5import} &
Raw data; user-defined metadata &
New HDF5 file &
Imports external data into HDF5, creating additional copies while potentially preserving metadata and semantics from the source format.\\
\hline

\texttt{h5jam}/\texttt{h5unjam} &
User block contents &
Extracted or embedded files &
Stores or extracts arbitrary information in the HDF5 user block, including scripts, XML documents, application metadata, or proprietary information.\\
\hline

\texttt{h5mkgrp} &
Structural metadata &
Modified HDF5 file &
Creates or modifies the group hierarchy. Limited direct privacy exposure but changes descriptive metadata and namespace organization.\\
\hline

\texttt{h5edit} &
Raw data; user-defined metadata &
Modified HDF5 file &
Creates or modifies attributes; modifies raw data. Limited direct privacy exposure but changes data and user-defined metadata.\\
\hline

\texttt{h5format\_convert} &
Raw data; structural metadata &
Converted HDF5 file &
Converts files between HDF5 format versions, creating additional copies while preserving or rewriting internal metadata structures.\\
\hline

\texttt{h5clear} &
Structural metadata &
Modified HDF5 file &
Clears SWMR status and metadata cache information. Primarily affects recovery metadata and file consistency information.\\

\end{longtable}
\endgroup

\paragraph{Note:} All HDF5 command-line tools are built on top of the HDF5 library 
and therefore inherit its capabilities. When VFDs or VOL connectors are configured, these 
tools may transparently access data stored in remote services, external storage systems, 
or non-HDF5 data sources exposed through connectors. Consequently, privacy exposure 
is not necessarily limited to the HDF5 file specified on the command line but may extend to 
the broader HDF5 ecosystem.
\subsection{Audit Findings and Mitigation Planning}

Formal core and official-tool findings are registered as
\code{CORE-SEC-001}, \code{CORE-SEC-002}, \code{CORE-TOOL-001},
\code{CORE-PRIV-001}, \code{CORE-SAFE-001}, and \code{CORE-SAFE-002} in
Appendix~\ref{app:findings-register}. The two shared-event security findings
are provisional historical parser-boundary scenarios; they do not assert that
every current HDF5 release contains every represented defect.
The assert-masking observation is separately registered as evidence record
\code{CORE-SEC-003} until exact source sites, supported release paths, and
consequences are established.

\paragraph{Priority reconciliation.}\mbox{}\\
The response order comes from Table~\ref{tab:findings-response-routing}, not
from the order of recommendations in this chapter. Confirmed affected scope for
\code{CORE-SEC-001}, \code{CORE-TOOL-001}, and \code{CORE-PRIV-001} receives
Near-term treatment; \code{CORE-SEC-002} receives Planned treatment. \code{CORE-SAFE-001} and
\code{CORE-SAFE-002} receive Near-term \emph{containment} --- capacity
headroom and monitoring, flush and checkpoint discipline, a rehearsed recovery
procedure, and independent copies of unregenerable data --- with their
structural treatment in \code{CORE-SAFE-S1} and \code{CORE-SAFE-S2} on the
Planned and Phase~3 horizons; \code{RM-2D} produces their acceptance evidence
rather than deciding whether they can be rated. The \code{CORE-SEC-003} catalog and
release-build differential study has a 30-day Phase~1 assurance target under
\code{RM-1E}; this scheduling choice does not assign it a risk tier. A
confirmed case is routed by its supported consequence and affected scope,
including under \code{CORE-SEC-001} for a memory-unsafe path or
\code{CORE-SEC-002} for termination or resource exhaustion. Phase~1 patching,
isolation, and privacy controls treat concrete populations. Phases~2--4 provide
reusable guardrails, class prevention, and sustainment; that structural work
may begin in parallel but does not extend an earlier response clock.

\subsubsection{Security and Safety}

The selected CVE history shows repeated malformed-file parsing pressure
against a complex binary metadata machine.
Affected variants should be patched or contained promptly. Prevention must not
stop at those local repairs: it should add systematic metadata validation at
decode boundaries, explicit invariants for each on-disk structure, fuzz corpora
tied to invariants, and hard separation between parser acceptance and internal
object construction.

\paragraph{Goal:} HDF5 should be able to parse any byte stream as hostile input
and either reject it cleanly or produce a validated internal object graph.
Anything else is security debt.

\paragraph{Structural prevention program:} maintain current-scope fixes or
containment while bootstrapping invariants, checked primitives, regression and
fuzz harnesses, review gates, and metrics in parallel. Migrate each decode
boundary in independently releasable tranches once its invariants and checked
primitives exist, and make the gates enforce each accepted tranche. This
program neither precedes nor delays Phase~1 treatment or Phase~2 release
safeguards.

The \code{CORE-S} identifiers below are stable identifiers for recommendations
in the core-security track. Their numbers do not express risk, urgency, effort,
dependency, or implementation order.

Parser normalization is a central structural control: decode file bytes into a
validated intermediate form before constructing internal objects so malformed
metadata cannot reach trusted C paths. It is a recurrence-prevention control,
not a substitute for patching or isolating a known affected release.

Our goal should be: \textit{make malformed HDF5 files boring}.
Every corrupted file should terminate in a controlled
\texttt{H5E\_BADVALUE}, \texttt{H5E\_OVERFLOW}, \texttt{H5E\_CORRUPT},
or \texttt{H5E\_RESOURCE} path. No segfaults. No assertions. No leaks.
No undefined behavior. No late-stage encode/flush crash from a bad
decode.

The current CVE list includes many issues in \texttt{H5FS},
\texttt{H5C}, \texttt{H5O}, \texttt{H5T}, \texttt{H5Z}, and
vector-copy paths, with repeated heap overflows, use-after-free,
null dereference, and resource-consumption outcomes.
HDF5’s own format documentation explains why this happens: a file is a
graph of groups, datasets, and named datatypes; on disk it is
represented through lower-level structures such as the superblock,
B-tree nodes, object headers, global heaps, local heaps, and
free-space metadata. That means the library is parsing a dense graph
of mutually referential, \textit{attacker-controlled} metadata.

\subsubsection{Prior Art: What Already Exists}
\label{sec:core-prior-art}

Several recommendations below describe controls that have already been
designed, and in some cases built and measured, in the H5Lens /
\code{hdf5-pickles} project~\cite{hdf5-pickles}. Presenting them as new proposals
would misstate both the effort involved and the state of the art inside the
programme, so this subsection records what exists before the recommendations
restate what is still missing.

H5Lens describes HDF5 on-disk metadata as machine-readable format definitions
and builds four things on top of them: \code{h5policy}, an independent
preflight oracle for untrusted files with named security profiles and a
regression corpus; a versioned invariant registry mapping each record family's
invariants to finding codes, enforcement points, tests, fuzz targets, and
migration status; \code{h5patch}, an evidence-gated repair planner; and a
differential measurement harness that records, per record family, what a
selected libhdf5 build actually did with a malformed specimen
\cite{h5policy,h5policy-invariant-registry,h5patch}. The accompanying design
document, \emph{A CVE Strategy for the HDF5 Library}, sets out the
decode-validate-construct architecture in more detail than this report
does~\cite{hdf5-cve-strategy,hdf5-bounded-raw-decode}.

\paragraph{Two places where the existing design is ahead of this report.}
This matters for sequencing, not just for credit.

\begin{enumerate}
    \item \textbf{Three validation scopes, not two phases.} \code{CORE-S2}
          below describes decode, validate, construct. The strategy document
          decomposes the middle step into \emph{record-local},
          \emph{aggregate-object}, and \emph{reference-graph} validation, which
          is the distinction that actually explains the harder CVE variants: a
          continuation message that references part of itself, a shared message
          that resolves through itself, or a compound member that extends past
          its enclosing datatype are all locally well-formed and fail only at a
          wider scope. A two-phase boundary does not by itself catch them.

    \item \textbf{Materialization versus activation.} The strategy document
          separates constructing a native object from performing operations
          with side effects --- cache publication, ID registration, plugin
          loading, decompression, user callbacks, and \emph{external-file
          access}. Its rule is that successfully parsing fields must not
          authorize any of those. This report has no equivalent concept, and
          the omission is consequential: that rule is the core-library-side
          statement of the same boundary failure this report rates as
          \code{APP-SEC-001} at the service boundary. The two framings are
          complementary, and a reader of this report alone would conclude that
          external-resource policy is purely a downstream concern.
\end{enumerate}

\paragraph{Re-scoped recommendations.}
Table~\ref{tab:core-prior-art-rescope} restates the affected recommendations
against what exists. The revised effort in the last column supersedes the
estimate carried in Appendix~\ref{app:recommendation-register} for these six
items, and the reuse is subject to the licensing boundary in
Section~\ref{sec:roadmap-licensing}: identifiers, boundaries, fixtures, and
measurements are consumable as data; GPLv3 implementation source is not
consumable into a BSD-licensed library.

\begingroup
\small
\setlength{\tabcolsep}{3pt}
\renewcommand{\arraystretch}{1.15}
\setlength{\LTleft}{-0.9in}
\setlength{\LTright}{-0.9in}
\begin{longtable}{@{}>{\raggedright\arraybackslash}p{0.11\widetablewidth}>{\raggedright\arraybackslash}p{0.34\widetablewidth}>{\raggedright\arraybackslash}p{0.38\widetablewidth}>{\raggedright\arraybackslash}p{0.15\widetablewidth}@{}}
\caption{Recommendations re-scoped against existing H5Lens work}
\label{tab:core-prior-art-rescope}\\
\hline
\textbf{Rec.} & \textbf{What already exists} &
\textbf{What this recommendation therefore is} & \textbf{Effort} \\
\hline
\endfirsthead

\caption[]{Recommendations re-scoped against existing H5Lens work, continued}\\
\hline
\textbf{Rec.} & \textbf{What already exists} &
\textbf{What this recommendation therefore is} & \textbf{Effort} \\
\hline
\endhead

\hline
\endfoot

\hline
\endlastfoot

\code{CORE-S1} &
A versioned, machine-readable invariant registry with per-record-family
invariants, finding codes, enforcement points, test and fuzz-target mappings,
and per-family migration status. &
\emph{Adopt the vocabulary; build the enforcement.} The design and enumeration
work is largely done. What remains is mapping each adopted invariant to an
always-active enforcement point \emph{inside libhdf5} and proving it in both
build modes. &
Large $\rightarrow$ \textbf{Medium} (design and enumeration reused;
enforcement is new) \\

\code{CORE-S2} &
A complete architectural design with three validation scopes, the
materialization/activation rule, a seven-phase incremental migration plan, and
a working bounded decoder demonstrating the contract. &
\emph{Adopt the design, including the two refinements above; migrate the
library.} The remaining work is the migration itself, which is genuinely large
and is not reduced by the existence of a design. &
Large $\rightarrow$ \textbf{Large} (unchanged; design reuse does not reduce
migration cost) \\

\code{CORE-S4} &
\code{h5policy} implements four profiles --- \code{legacy},
\code{trusted-fast}, \code{untrusted-strict}, \code{forensic} --- with the same
names and substantially the same semantics this recommendation proposes,
including feature policy and resource budgets. &
\emph{Port a proven profile design into the library API.} The semantics have
been exercised against a corpus; the open work is the \code{H5P} surface,
default-state behavior, and enforcement inside libhdf5 rather than beside it. &
Large $\rightarrow$ \textbf{Medium} \\

\code{CORE-S5} &
A mutation harness with per-family recipes, truncation sweeps, and a
boundary-value inventory identifying which values the corpus lacks. Its own
coverage report records that OSS-Fuzz integration is absent for every family. &
\emph{Adopt the harness design and close the two named gaps:} OSS-Fuzz
integration and per-family fuzz targets, which exist for only a small minority
of families. This recommendation should stop proposing OSS-Fuzz as a new idea
and start treating it as a known, scheduled gap. &
Large $\rightarrow$ \textbf{Medium} \\

\code{CORE-S6} &
A malformed-input corpus with per-fixture expected-outcome records, a valid
control set, and a manifest --- structurally the \code{expected.yml} design
sketched below. &
\emph{Adopt the corpus and its expectation format;} add HDF5-side CI
integration and the debug/\code{-DNDEBUG} parity requirement. Subject to
per-specimen provenance review before redistribution. &
Medium $\rightarrow$ \textbf{Small} \\

\code{APP-S2} &
\code{h5policy} \emph{is} a working preflight validator for untrusted files,
invocable per profile and emitting structured results. &
\emph{Two distinct deliverables, previously conflated.} A validator exists and
can be used today by deployments willing to run an external GPLv3 tool. What
does not exist is the enforcement coupling this recommendation requires ---
binding the preflight verdict to the same open handle or verified file
identity used at load time --- which is the part that makes preflight an
authorization rather than an advisory. &
Large $\rightarrow$ \textbf{Medium} (validator reused; coupling is the real
work) \\

\end{longtable}
\endgroup

\paragraph{What this report adds that the existing work does not cover.}
Stating the reuse honestly also requires stating the complement. The
application-boundary findings, the independent-implementation and
interoperability analysis, the separation of inherent from residual risk with
response clocks and accountable owners, and the attribution of recommendations
to enforcing actors across downstream projects have no counterpart in H5Lens,
which is deliberately a format-and-oracle project rather than an
ecosystem-governance one. The two are complementary. The defect corrected here
is that this report previously described the overlap as future work.

\paragraph{An unclosed loop.}
The existing case records report that a substantial majority of their tracked
divergences are \emph{oracle-hardened}: the oracle rejects the specimen and the
libhdf5 side is unfixed. Two of those concern external-file-list handling, and
one --- an external data file opened before the wrapped segment is validated
--- is recorded as
uninvestigated~\cite{h5policy-case-efl-open}. If substantiated, that is a
core-library instance of the activation-before-validation rule above and bears
directly on \code{APP-SEC-001}. This report did not evaluate it, and
\code{RM-1E} should take the full set of oracle-hardened cases as triage input.

\subsubsection*{Recommendation CORE-S0: Separate defect closure from recurring-class closure}

Adopt the following rule:

\textit{A verified local repair may close the affected defect, but it must not
be represented as elimination of the recurring class without structural
evidence.}

For every new CVE or crash:

\begingroup
\small
\setlength{\tabcolsep}{3pt}
\renewcommand{\arraystretch}{1.15}

\begin{longtable}{@{}p{0.30\linewidth}p{0.66\linewidth}@{}}
\caption{Required artifacts for defect closure and recurring-class tracking}
\label{tab:cve-closure-artifacts}\\
\hline
\textbf{Required artifact} & \textbf{Purpose} \\
\hline
\endfirsthead

\caption[]{Required artifacts for defect closure and recurring-class tracking, continued}\\
\hline
\textbf{Required artifact} & \textbf{Purpose} \\
\hline
\endhead

\hline
\endfoot

\hline
\endlastfoot

crashing input & becomes a permanent corpus seed \\
minimized input & used in quick PR regression \\
root-cause category & maps to a known invariant class \\
invariant update & records and advances prevention of sibling bugs \\
sanitizer regression & proves no crash, no leak, no UAF \\
negative test & checks the same controlled failure in assertion-enabled and
\code{-DNDEBUG} builds \\
reviewer checklist update & prevents reintroduction \\
\end{longtable}
\endgroup

Variant closure requires evidence that the affected path is repaired or
contained and that its negative and sanitizer regressions pass. A local
\texttt{if} may therefore be a valid treatment, including for an embargoed
fix, while the corresponding \textbf{class-prevention record} remains open
until the invariant framework and related gates cover sibling cases.

Recent commits show the right direction, but too locally. One fix
added a sanity check immediately after reading file-space page size
from the file and used the same fuzzer to verify that the previous
assertion became a corrupted-file error.~\cite{commit_804f3ba}
Another hardened \texttt{H5C\_\_reconstruct\_cache\_entry()} by adding
buffer-size tracking, overflow checks, input validation, and safe cleanup.~\cite{commit_3914bb7}
Those are exactly the moves to systematize.

\subsubsection*{Recommendation CORE-S1: Build an invariant registry for the file format}

\noindent\emph{Prior art:} see Table~\ref{tab:core-prior-art-rescope}; this item is re-scoped against existing work.
Create a versioned document and machine-readable metadata file, for example:

\begin{verbatim}
security/invariants/
  superblock.yml
  btree_v2.yml
  chunk_records.yml
  extensible_array.yml
  object_header.yml
  dataspace.yml
  datatype.yml
  layout.yml
  filters.yml
  heaps.yml
  fractal_heap.yml
  free_space.yml
  links.yml
  shared_object_header_messages.yml
  cache_image.yml
\end{verbatim}

Each entry should define:

\begin{minted}{yaml}
id: H5O-CONT-001
component: object_header
structure: continuation_message
condition: continuation_chunk_size > 0
failure: H5E_CORRUPT
checked_at:
  - H5O__cont_decode
  - H5O__chunk_deserialize
release_enforced: true
diagnostic_assert_sites: []
fuzz_targets:
  - fuzz_object_header
cves:
  - CVE-2025-6856
  - CVE-2025-6816
\end{minted}

The following are candidate initial invariant-registry families. Tranche order
must be recorded in the implementation plan from affected scope, coverage gaps,
and concrete enabling dependencies; this list is not a risk rank.

\begingroup
\small
\setlength{\tabcolsep}{3pt}
\renewcommand{\arraystretch}{1.15}

\setlength{\LTleft}{-0.7in}
\setlength{\LTright}{-0.7in}

\begin{longtable}{@{}p{0.16\widetablewidth}p{0.80\widetablewidth}@{}}
\caption{Candidate initial invariant-registry families}
\label{tab:invariant-families-first}\\
\hline
\textbf{Area} & \textbf{Candidate invariants} \\
\hline
\endfirsthead

\caption[]{Candidate initial invariant-registry families, continued}\\
\hline
\textbf{Area} & \textbf{Candidate invariants} \\
\hline
\endhead

\hline
\endfoot

\hline
\endlastfoot

superblock &
offset size and length size are valid; base address, EOF address, driver info
address, and root object address are inside the file or explicitly undefined;
EOF does not wrap \\
object headers &
message length fits inside chunk; continuation chunks are nonzero; continuations
do not self-reference; continuations are acyclic; message flags match message
type; unknown messages cannot be cast into known message structures \\
free-space metadata &
page size is nonzero and below max; section size fits page/file bounds;
serialized section count matches decoded section count \\
heaps &
heap object offset and size are inside heap; free-list blocks do not overlap;
free-block list is acyclic; no allocation uses unchecked \code{count * size} \\
dataspaces &
rank is bounded; dimension product is checked; selected point count is checked;
hyperslab stride, count, block, and start do not overflow \\
datatypes &
storage size is nonzero; precision fits storage size; member offsets fit
compound size; string length is bounded; vlen/ref types cannot escape validation \\
layouts &
compact data size equals \code{dataspace elements * datatype size}; chunk
dimensions are nonzero; chunk product is checked; chunk index references are
inside file bounds \\
filters &
parameter count is bounded; decompressed size is bounded; scale-offset precision
is validated; filter output ratio cannot exceed configured resource limits \\
cache image &
cache-entry count, address, size, and alignment are bounded; decoding consumes
exactly the expected number of bytes; cleanup is safe after partial decode \\
tools &
string rendering is length-bounded; conversion tools never trust input
dimensions; tool-only parsers use the same checked arithmetic wrappers as the
library \\
\end{longtable}
\endgroup

The seven component families named in issue~\#87 are the mandatory first
audit perimeter for assert-masked checks. That perimeter is not itself a risk
rank. Within it, first migrate a candidate that is a sole release-reachable
guard for memory bounds, pointer arithmetic, allocation, copy, or traversal;
then address termination-only paths and redundant diagnostics. Record the
reachability and consequence evidence so a raw candidate count cannot
substitute for prioritization.

This registry becomes the source of truth for code review, tests, fuzzing, and documentation.

\subsubsection*{Recommendation CORE-S2: Introduce a two-phase parser boundary}

\noindent\emph{Prior art:} see Table~\ref{tab:core-prior-art-rescope}; this item is re-scoped against existing work.
Right now, too many paths appear to do this:
\begin{verbatim}
file bytes -> native internal object
           -> later code assumes validity
\end{verbatim}

Change the architecture to:
\begin{verbatim}
file bytes -> bounded raw decode
           -> semantic validation
           -> native object construction
\end{verbatim}

Concrete rule: \textit{No public or internal runtime object is constructed from
file metadata until the corresponding validator has passed.}

Implement the pattern by component:

\begin{minted}{c}
herr_t H5O_msg_decode_raw(const uint8_t *buf, size_t buf_size, H5O_msg_raw_t *raw);
herr_t H5O_msg_validate(const H5F_t *f, const H5O_msg_raw_t *raw, H5O_validation_ctx_t *ctx);
herr_t H5O_msg_construct(const H5O_msg_raw_t *raw, H5O_msg_t *out);    
\end{minted}

Do not let decode functions allocate complex structures before basic
bounds are known. A raw decoded object should hold values, offsets,
and spans, not trusted pointers into long-lived internal state.

All decode functions that consume file bytes should take one of these forms:
\begin{minted}{c}
herr_t decode(const H5F_t *f, const uint8_t *buf, size_t buf_size, out_t *out);

herr_t decode_cursor(H5_decode_cursor_t *cur, out_t *out);
\end{minted}

Ban this pattern in security-sensitive decode code:
\begin{minted}{c}
const uint8_t **pp
\end{minted}
unless there is also an explicit \texttt{p\_end} or remaining-size counter. The cache-image fix already moved in this direction
by adding a buffer-size argument and bounds checks.~\cite{commit_3914bb7}

\subsubsection*{Recommendation CORE-S3: Centralize checked arithmetic and allocation}

Most parser bugs become exploitable when metadata reaches allocation, copy, or pointer arithmetic. Add a small, mandatory checked-arithmetic layer.

Example API:
\begin{minted}{c}
herr_t H5_checked_add_hsize(hsize_t a, hsize_t b, hsize_t *out);
herr_t H5_checked_mul_hsize(hsize_t a, hsize_t b, hsize_t *out);
herr_t H5_checked_addr_add(haddr_t base, hsize_t len, haddr_t eoa, haddr_t *end);
herr_t H5_checked_alloc_size(size_t count, size_t elem_size, size_t *bytes);
\end{minted}

Then make these rules review-blocking:

\begingroup
\small
\setlength{\tabcolsep}{3pt}
\renewcommand{\arraystretch}{1.15}

\setlength{\LTleft}{-0.9in}
\setlength{\LTright}{-0.9in}

\begin{longtable}{@{}p{0.48\widetablewidth}p{0.48\widetablewidth}@{}}
\caption{Review-blocking parser arithmetic and allocation patterns}
\label{tab:review-blocking-parser-patterns}\\
\hline
\textbf{Ban} & \textbf{Replacement} \\
\hline
\endfirsthead

\caption[]{Review-blocking parser arithmetic and allocation patterns, continued}\\
\hline
\textbf{Ban} & \textbf{Replacement} \\
\hline
\endhead

\hline
\endfoot

\hline
\endlastfoot

\code{malloc(n * size)} & \code{H5MM_calloc_checked(n, size)} \\
\code{addr + size} & \code{H5_checked_addr_add(addr, size, eoa, \&end)} \\
\code{memcpy(dst, src, n)} from file-derived size &
\code{H5_checked_memcpy(dst, dst_size, src, src_size, n)} \\
pointer increment without \code{p_end} & cursor-based decode \\
raw recursion through object graph & bounded traversal context \\
\end{longtable}
\endgroup

The 2024 security commit added multiplication-overflow checks for
dataset storage size and checks against file bounds.~\cite{commit_ce53bc0}
That pattern should become a library-wide rule, not a local repair.

\subsubsection*{Recommendation CORE-S4: Create security profiles for file opening}

\noindent\emph{Prior art:} see Table~\ref{tab:core-prior-art-rescope}; this item is re-scoped against existing work.
Add explicit modes. Do not rely on users to infer safe behavior.
\begin{minted}{c}
typedef enum {
    H5_SECURITY_PROFILE_LEGACY,
    H5_SECURITY_PROFILE_TRUSTED_FAST,
    H5_SECURITY_PROFILE_UNTRUSTED_STRICT,
    H5_SECURITY_PROFILE_FORENSIC
} H5_security_profile_t;
\end{minted}

Suggested semantics:

\begingroup
\small
\setlength{\tabcolsep}{3pt}
\renewcommand{\arraystretch}{1.15}

\begin{longtable}{@{}p{0.24\linewidth}p{0.72\linewidth}@{}}
\caption{Security profile semantics}
\label{tab:security-profile-semantics}\\
\hline
\textbf{Profile} & \textbf{Behavior} \\
\hline
\endfirsthead

\caption[]{Security profile semantics, continued}\\
\hline
\textbf{Profile} & \textbf{Behavior} \\
\hline
\endhead

\hline
\endfoot

\hline
\endlastfoot

\code{legacy} & current compatibility bias \\
\code{trusted_fast} & normal validation, generous resource limits \\
\code{untrusted_strict} & strict metadata validation, plugin loading disabled unless allowlisted, resource limits enforced \\
\code{forensic} & never repairs silently, reports all structural anomalies, best-effort inspection without unsafe traversal \\
\end{longtable}
\endgroup

Expose it through file-access property lists:
\begin{minted}{c}
H5Pset_security_profile(fapl, H5_SECURITY_PROFILE_UNTRUSTED_STRICT);
H5Pset_max_metadata_bytes(fapl, 256 * 1024 * 1024);
H5Pset_max_object_header_chunks(fapl, 1024);
H5Pset_max_btree_depth(fapl, 64);
H5Pset_max_filter_expansion_ratio(fapl, 100);
H5Pset_plugin_policy(fapl, H5_PLUGIN_POLICY_DISABLED);
\end{minted}

This matters because HDF5 supports dynamically loaded plugins, and the
runtime API can enable or disable plugin classes; the docs also state
that \texttt{HDF5\_PLUGIN\_PRELOAD=::} disables all dynamic plugins.
Turn that capability into a first-class security policy, not an
environment-variable trick.

\subsubsection*{Recommendation CORE-S5: Fuzz by structure, not just by whole-file crashes}

\noindent\emph{Prior art:} see Table~\ref{tab:core-prior-art-rescope}; this item is re-scoped against existing work.
Whole-file fuzzing is necessary but insufficient. It finds symptoms. Structure-aware fuzzing finds families.

Treat the following as an initial harness-investment portfolio, not a finding
rank or response order. Known finding fixtures and changed paths enter
regression immediately; new harnesses may proceed in parallel, with tranche
order recorded from shared-enablement and coverage evidence.

\begingroup
\small
\setlength{\tabcolsep}{3pt}
\renewcommand{\arraystretch}{1.15}

\setlength{\LTleft}{-0.9in}
\setlength{\LTright}{-0.9in}

\begin{longtable}{@{}p{0.38\widetablewidth}p{0.58\widetablewidth}@{}}
\caption{Candidate structure-aware fuzzing portfolio}
\label{tab:structure-aware-fuzzing-sequence}\\
\hline
\textbf{Candidate target} & \textbf{Coverage rationale} \\
\hline
\endfirsthead

\caption[]{Candidate structure-aware fuzzing portfolio, continued}\\
\hline
\textbf{Candidate target} & \textbf{Coverage rationale} \\
\hline
\endhead

\hline
\endfoot

\hline
\endlastfoot

object-header message decoder & shared decode paths and continuation structures \\
datatype decoder and converter & complex nested metadata and conversion operations \\
dataspace and selection decoder & rank, count, hyperslab, and point arithmetic \\
layout and chunk index decoder & storage-size, address, and chunk-product checks \\
free-space and local/global heap decoders & historical overflow and list-corruption variants \\
filter pipeline and scale-offset filter & compression and decompression boundaries \\
metadata cache image reconstruction & complex bounds, lifecycle, and cleanup paths \\
\code{h5dump}/\code{h5repack} rendering paths & tool paths that add handling and rendering behavior \\
\end{longtable}
\endgroup

Use three fuzzing layers:

\begingroup
\small
\setlength{\tabcolsep}{3pt}
\renewcommand{\arraystretch}{1.15}

\setlength{\LTleft}{-0.9in}
\setlength{\LTright}{-0.9in}

\begin{longtable}{@{}p{0.28\widetablewidth}p{0.68\widetablewidth}@{}}
\caption{Fuzzing layers}
\label{tab:fuzzing-layers}\\
\hline
\textbf{Layer} & \textbf{Description} \\
\hline
\endfirsthead

\caption[]{Fuzzing layers, continued}\\
\hline
\textbf{Layer} & \textbf{Description} \\
\hline
\endhead

\hline
\endfoot

\hline
\endlastfoot

byte-level fuzzing & current randomized mutation against APIs \\
structure-aware mutation & valid superblock with corrupted object headers, valid object header with bad continuation, valid datatype with bad member offset \\
invariant-guided generation & generate files that violate exactly one invariant at a time \\
\end{longtable}
\endgroup

OSS-Fuzz already supports libFuzzer, AFL++, Honggfuzz, Centipede, sanitizers,
and distributed execution for C/C++ projects. Use OSS-Fuzz for long-running
public coverage and ClusterFuzzLite or local CI for PR-level checks.

\begingroup
\small
\setlength{\tabcolsep}{3pt}
\renewcommand{\arraystretch}{1.15}

\begin{longtable}{@{}p{0.24\linewidth}p{0.72\linewidth}@{}}
\caption{Fuzzing acceptance criteria}
\label{tab:fuzzing-acceptance-criteria}\\
\hline
\textbf{Stage} & \textbf{Gate} \\
\hline
\endfirsthead

\caption[]{Fuzzing acceptance criteria, continued}\\
\hline
\textbf{Stage} & \textbf{Gate} \\
\hline
\endhead

\hline
\endfoot

\hline
\endlastfoot

PR & run minimized CVE corpus under \code{ASan}/\code{UBSan}/\code{LSan} \\
nightly & run structure fuzzers for core parsers \\
weekly & report coverage by component, not just total line coverage \\
release candidate & 72-hour fuzz soak on changed parser areas \\
CVE closure & crashing input added to corpus and mapped to invariant \\
\end{longtable}
\endgroup

\subsubsection*{Recommendation CORE-S6: Make malformed-input regression cheap}

\noindent\emph{Prior art:} see Table~\ref{tab:core-prior-art-rescope}; this item is re-scoped against existing work.
Create this repository layout:
\begin{verbatim}
test/security/
  corpus/
    cve/
      CVE-2025-7069/
        input.h5
        expected.yml
      CVE-2025-6269/
        input.h5
        expected.yml
    invariants/
      H5O-CONT-001/
      H5T-SIZE-002/
  harnesses/
    fuzz_object_header.c
    fuzz_datatype.c
    fuzz_dataspace.c
    fuzz_layout.c
    fuzz_filters.c
  expected_errors/
    corrupt_file.yml
\end{verbatim}

Each expected.yml should include:
\begin{minted}{yaml}
expected:
  exit: controlled_error
  allowed_errors:
    - H5E_CORRUPT
    - H5E_BADVALUE
    - H5E_OVERFLOW
forbidden:
  - SIGSEGV
  - SIGABRT
  - leak
  - use_after_free
  - timeout
invariant:
  - H5FS-PAGE-001
\end{minted}

This gives maintainers a fast way to prove that old crashes remain dead.

\subsubsection*{Recommendation CORE-S7: Add a security-review checklist that blocks risky parser changes}

Every PR touching decode, encode, object lifecycle, filters, tools, or plugin
loading should answer:

\begingroup
\small
\setlength{\tabcolsep}{3pt}
\renewcommand{\arraystretch}{1.15}

\setlength{\LTleft}{-0.9in}
\setlength{\LTright}{-0.9in}

\begin{longtable}{@{}p{0.58\widetablewidth}p{0.38\widetablewidth}@{}}
\caption{Security-review checklist for parser-adjacent changes}
\label{tab:security-review-checklist}\\
\hline
\textbf{Question} & \textbf{Required answer} \\
\hline
\endfirsthead

\caption[]{Security-review checklist for parser-adjacent changes, continued}\\
\hline
\textbf{Question} & \textbf{Required answer} \\
\hline
\endhead

\hline
\endfoot

\hline
\endlastfoot

Does this read file-derived bytes? & list buffer and bounds source \\
Does this allocate using file-derived values? & show checked multiplication \\
Does this copy memory using file-derived sizes? & show source and destination bounds \\
Does this traverse links, chunks, heaps, B-trees, or continuations? & show cycle and depth bounds \\
Does this construct native objects before validation? & explain why or refactor \\
Does any file-derived validity predicate appear only in an assertion? & show
the always-active check and matching \code{-DNDEBUG} negative test \\
Does this add or change cleanup paths? & show partial-init safety \\
Does this invoke plugins or filters? & show policy interaction \\
Does this alter error handling? & show malformed-input test \\
\end{longtable}
\endgroup

Add automated enforcement where possible:

\begin{verbatim}
fail PR if new decode function lacks size parameter
fail PR if malloc(count * size) appears outside checked wrapper
fail PR if memcpy appears in parser code without checked wrapper annotation
fail PR if new parser file lacks fuzz target registration
fail PR if new object traversal lacks max-depth or visited-set logic
fail PR if file-derived validity is assert-only
run malformed corpus in debug and -DNDEBUG builds
require the same controlled rejection in both
\end{verbatim}

This will annoy contributors at first. Good. That friction is cheaper than post-release CVEs.

\subsubsection*{Recommendation CORE-S8: Separate compatibility from safety}

HDF5’s backward compatibility is valuable, but it creates a security trap: old,
permissive parsing keeps malformed legacy structures alive.

Implement this policy:

\begingroup
\small
\setlength{\tabcolsep}{3pt}
\renewcommand{\arraystretch}{1.15}

\setlength{\LTleft}{-0.7in}
\setlength{\LTright}{-0.7in}

\begin{longtable}{@{}p{0.42\widetablewidth}p{0.54\widetablewidth}@{}}
\caption{Compatibility and safety profile policy}
\label{tab:compatibility-safety-profile-policy}\\
\hline
\textbf{Case} & \textbf{Behavior} \\
\hline
\endfirsthead

\caption[]{Compatibility and safety profile policy, continued}\\
\hline
\textbf{Case} & \textbf{Behavior} \\
\hline
\endhead

\hline
\endfoot

\hline
\endlastfoot

valid legacy file & accepted \\
ambiguous legacy construct & accepted only in \code{legacy} or \code{trusted_fast}; warning in \code{forensic}; rejected in \code{untrusted_strict} \\
structurally invalid file that old versions tolerated & rejected in all non-legacy profiles \\
parser repair needed & repair only through explicit tool, never silently during normal open \\
unknown message with unsafe flags & rejected unless profile explicitly allows preservation \\
\end{longtable}
\endgroup

The point is not to break legitimate archives. The point is to stop treating
corrupted metadata as compatibility.

\subsubsection*{Recommendation CORE-S9: Harden tools as separate attack surfaces}

Do not assume \code{h5dump}, \code{h5repack}, converters, and test utilities are
less important. They are often the first thing users run on untrusted files.

Actions:

\begingroup
\small
\setlength{\tabcolsep}{3pt}
\renewcommand{\arraystretch}{1.15}

\begin{longtable}{@{}p{0.30\linewidth}p{0.66\linewidth}@{}}
\caption{Tool hardening actions}
\label{tab:tool-hardening-actions}\\
\hline
\textbf{Tool class} & \textbf{Change} \\
\hline
\endfirsthead

\caption[]{Tool hardening actions, continued}\\
\hline
\textbf{Tool class} & \textbf{Change} \\
\hline
\endhead

\hline
\endfoot

\hline
\endlastfoot

inspection tools & use \code{H5_SECURITY_PROFILE_FORENSIC} by default \\
conversion tools & use \code{H5_SECURITY_PROFILE_UNTRUSTED_STRICT} by default \\
rendering paths & bounded strings only \\
legacy converters & disable by default; build behind explicit \code{HDF5_ENABLE_LEGACY_UNSAFE_TOOLS=ON} \\
plugin-dependent tools & warn or require allowlist \\
\end{longtable}
\endgroup

A dangerous file should not become more dangerous because the user tried to
inspect it.

\subsubsection*{Recommendation CORE-S10: Plugin policy and signed extensions}

Plugins are executable extension code requiring a distinct control domain; this
statement is not a risk rating. They are loaded at runtime.
HDF5’s plugin system supports dynamically loaded filter plugins and lets users
implement custom filters. That is useful, but it must not be invisible in
untrusted mode.

Implement:
\begin{verbatim}
default profile:
  trusted_fast: existing plugin behavior
  untrusted_strict: plugins disabled unless allowlisted
  forensic: plugins disabled; report missing filter
\end{verbatim}

Add a plugin manifest:
\begin{minted}{yaml}
plugin_id: 32008
name: blosc
version: 1.0.0
sha256: ...
signer: ...
allowed_profiles:
  - trusted_fast
  - untrusted_strict
capabilities:
  - filter_decode
limits:
  max_expansion_ratio: 100
  max_alloc_bytes: 1073741824
\end{minted}

For affected deployments, disable plugins or restrict loading to protected,
explicitly approved paths and artifacts within the applicable response clock.
Hash allowlists may be an interim control; supported signature verification and
resource budgets should be configured and tested independently when applicable.
Neither waits for a complete PKI nor extends a finding's response horizon, and
a valid signature does not make a buggy plugin resource-safe.

\subsubsection*{Recommendation CORE-S11: Coordinate supply-chain controls with parser hardening}

Supply-chain work can proceed in parallel with parser hardening. Assign it an
owner and scope independently, and record any concrete cross-track dependency
rather than inferring order from an unscored comparison of the risk areas.

Minimum controls:

\begingroup
\small
\setlength{\tabcolsep}{3pt}
\renewcommand{\arraystretch}{1.15}

\begin{longtable}{@{}p{0.32\linewidth}p{0.64\linewidth}@{}}
\caption{Minimum supply-chain controls}
\label{tab:minimum-supply-chain-controls}\\
\hline
\textbf{Control} & \textbf{Target} \\
\hline
\endfirsthead

\caption[]{Minimum supply-chain controls, continued}\\
\hline
\textbf{Control} & \textbf{Target} \\
\hline
\endhead

\hline
\endfoot

\hline
\endlastfoot

OpenSSF Scorecard & weekly scheduled run, alerts in security tab \\
CodeQL & required for PRs touching C/C++ parser surfaces \\
pinned GitHub Actions & no floating third-party action refs \\
SBOM & release artifacts include source and binary SBOMs \\
signed release artifacts & signatures plus checksums \\
SLSA provenance & release builds produce verifiable provenance \\
\end{longtable}
\endgroup

OpenSSF Scorecard has an official GitHub Action, and public repositories can use
it without cost. The SLSA GitHub generator provides tooling for SLSA provenance
on GitHub Actions, including provenance that users can verify to understand how
an artifact was built.

\subsubsection*{Recommendation CORE-S12: Metrics that force the right behavior}

Use metrics that measure class elimination, not activity.

\begingroup
\small
\setlength{\tabcolsep}{3pt}
\renewcommand{\arraystretch}{1.15}

\begin{longtable}{@{}p{0.70\linewidth}p{0.26\linewidth}@{}}
\caption{Security program metrics}
\label{tab:security-program-metrics}\\
\hline
\textbf{Metric} & \textbf{Target} \\
\hline
\endfirsthead

\caption[]{Security program metrics, continued}\\
\hline
\textbf{Metric} & \textbf{Target} \\
\hline
\endhead

\hline
\endfoot

\hline
\endlastfoot

CVE corpus pass rate under ASan/UBSan/LSan & 100\% \\
parser decode functions with explicit buffer size or cursor & 100\% \\
parser allocations through checked wrappers & 100\% \\
parser memcpy/memmove through checked wrappers or reviewed exemption & 100\% \\
structures in the versioned declared tranche with invariant files & 100\% \\
structures in the versioned declared tranche with registered fuzz targets & 100\% \\
cataloged file-derived assertions classified for dominance, reachability, and consequence & 100\% \\
declared file-derived invariants enforced in \code{-DNDEBUG} builds & 100\% \\
new CVEs that add invariant updates & 100\% \\
time from new crash to minimized corpus seed & under 48 hours \\
time from minimized seed to invariant classification & under 5 business days \\
sanitizer-clean nightly fuzz runs & sustained \\
untrusted profile test coverage & every release \\
\end{longtable}
\endgroup

Track the versioned tranche denominator and results in a public dashboard. The
dashboard does not fix defects; it makes declared coverage and gaps auditable.

\subsubsection{Safety: Durability and Recovery}

The security track above answers ``a hostile file arrived.'' The two
recommendations here answer ``my file did not survive.'' They address
\code{CORE-SAFE-001} and \code{CORE-SAFE-002}, and neither has a counterpart
elsewhere in this report. The \code{CORE-SAFE-S} identifiers name
recommendations only; their numbers express no risk, urgency, effort,
dependency, or implementation order.

Both findings are routed Near-term for \emph{containment} and Planned or later
for \emph{structural treatment}. The containment obligations are operational
and available today: capacity headroom with free-space monitoring, flush and
checkpoint discipline, a rehearsed recovery procedure, and an independent copy
of anything that cannot be regenerated. Nothing below is a prerequisite for
those.

\subsubsection*{Recommendation CORE-SAFE-S1: Give HDF5 an \code{ENOSPC} story}

Today an application can receive an HDF5 write, flush, or close failure, but it
has no structured callback identifying the VFD's \code{ENOSPC} event, no
supported protocol for choosing retry versus abandonment at that boundary, and
no library state that marks the resulting file as untrustworthy. Four candidate
mechanisms have already been proposed in the public design discussion; they are
complementary rather than alternatives, and they differ sharply in
cost~\cite{forum-enospc-11714}.

\begingroup
\small
\setlength{\tabcolsep}{3pt}
\renewcommand{\arraystretch}{1.15}
\begin{longtable}{@{}>{\raggedright\arraybackslash}p{0.26\linewidth}>{\raggedright\arraybackslash}p{0.44\linewidth}>{\raggedright\arraybackslash}p{0.24\linewidth}@{}}
\caption{Candidate \code{ENOSPC} mechanisms}
\label{tab:enospc-mechanisms}\\
\hline
\textbf{Mechanism} & \textbf{What it provides} & \textbf{Cost and caveat} \\
\hline
\endfirsthead
\caption[]{Candidate \code{ENOSPC} mechanisms, continued}\\
\hline
\textbf{Mechanism} & \textbf{What it provides} & \textbf{Cost and caveat} \\
\hline
\endhead
\hline
\endfoot
\hline
\endlastfoot

VFD-level error callback &
The caller is told the driver hit \code{ENOSPC} and chooses retry, fail fast,
or panic. This is the load-bearing item: it converts a silent unsafe state into
an application decision. &
Small to medium. Deliverable as a pass-through VFD wrapper first, without a
public API commitment. \\

Poisoned-file state &
A file that suffered a write failure is durably marked suspect, so a later open
cannot silently treat it as sound. &
Medium; needs a format-visible flag and a defined clearing procedure. Pairs
with \code{h5clear}. \\

Space reservation (\code{posix\_fallocate}) &
Capacity is claimed before the write, converting a mid-write failure into an
up-front one. &
Medium; filesystem-dependent, and not available or meaningful on every backing
store. \\

Pre-flush capacity check (\code{statvfs}) &
A watchdog refuses to begin a flush that plainly cannot complete. &
Small, but advisory only --- it narrows the window and does not close it. A
check that passes does not guarantee the flush succeeds. \\

\end{longtable}
\endgroup

Sequence them by that cost profile: the callback and the pre-flush check first,
because they are small and independently useful; reservation and the poisoned
state after, because they touch the format and the VFD contract. The public
design discussion proposes a pass-through VFD wrapper as a low-friction proving
ground before any public API change, and this report endorses that route.

\noindent\textbf{Acceptance evidence.}
Fault injection returns \code{ENOSPC} at each of write, flush, and close, for
each supported VFD and buffering profile. For every injected point: the
application receives a distinguishable \code{ENOSPC}-class error rather than a
generic failure or a success; the last durable state is recorded; a subsequent
open either succeeds, fails with a diagnostic naming the condition, or reports
the file as suspect --- never silently returns wrong data; and the documented
recovery procedure is executed and its result recorded. A run in which the
library reports success and the file is later found inconsistent is a failure
of this recommendation regardless of the error codes returned.

\subsubsection*{Recommendation CORE-SAFE-S2: State, then improve, the crash-consistency guarantee}

The immediate deliverable is not a mechanism. It is a written statement of what
HDF5 guarantees today when a writer dies mid-update --- for each of the default
single-file layout, SWMR, and the split and multi drivers --- and, explicitly,
what it does not. Practitioners currently infer this from forum threads and
folklore, which is why the same questions recur across sixteen
years~\cite{forum-crash-state-1598,forum-flush-crash-3481,
forum-crash-corruption-7581}. A documented weak guarantee is more useful than
an undocumented one, because it lets an application decide whether it needs
checkpointing, an external journal, or a different store.

The structural work is a multi-release programme and is already under
investigation: metadata journaling, write-ahead logging, mechanisms derived
from SWMR, and checkpointing~\cite{forum-crashproofing-12003}. This report does
not select among them. It records three constraints on whichever is chosen.

\begin{enumerate}
    \item \textbf{Recovery tooling is part of the guarantee, not an
          afterthought.} Whether a file reopens after \code{h5clear} is the
          difference between I2 and I3 in \code{CORE-SAFE-002}. Improving
          crash consistency without improving diagnosis and repair moves files
          from one failure class to another rather than reducing harm.
    \item \textbf{Silent wrong data is the outcome to design against.} A file
          that fails to open is a bounded operational problem. A file that
          opens and returns data from an inconsistent state is a scientific
          integrity problem, and it is the one that will not be noticed.
    \item \textbf{The guarantee must be testable by fault injection,} and its
          test must run in the release build configuration, not only under
          assertions --- the same discipline \code{CORE-S7} applies to parser
          invariants.
\end{enumerate}

\noindent\textbf{Acceptance evidence.}
Published documentation states the crash-consistency guarantee per layout and
driver, including its limits. A fault-injection harness kills the writer at a
matrix of points across metadata update, flush, and close, and for every point
records whether the file opens, opens after repair, or does not open, and
whether any data read back differs from what was acknowledged as written. The
last case is reported separately and treated as a defect, never as expected
behaviour. Results are published per supported release so that the guarantee
can be seen to hold or to change.

\subsubsection{Privacy}

This section summarizes the Privacy Exposure Model's mechanisms and the bounded
\code{CORE-PRIV-001} finding; the other mechanisms are not separately rated.
It then describes corresponding mitigation strategies.

\paragraph{Overview of Privacy Threat Categories}\mbox{}\\

The HDF5 Privacy Exposure Model ~\cite{hdf5-privacy-exposure-model} identifies 
seven exposure families that describe mechanisms through which sensitive
information may be unintentionally exposed in HDF5. These mechanisms can arise
during normal scientific data creation, processing, sharing, and preservation;
their frequency, affected population, and consequence remain
deployment-specific.

Several threat categories involve direct disclosure of information stored within HDF5 files. 
Semantic Metadata Exposure (P1) occurs when descriptive information, such as object names, attributes,
datatype member names, or other user-defined metadata, reveals sensitive contextual 
information even when the underlying scientific data are protected. Structural or On-disk
Leakage (P2) arises from examination of HDF5's internal storage structures,
including file layout, storage organization, and implementation metadata, which may 
disclose characteristics of the stored data without requiring HDF5 software. 
Raw-data Pattern Leakage (P3) addresses situations in which information about scientific data
becomes observable through storage patterns, embedded textual descriptions, or 
other implementation artifacts rather than through intended data access.

Another group of threats results from routine scientific data processing. Unintended
Inclusion During Processing (P4) occurs when software automatically generates, expands,
or preserves additional information during data conversion or processing.
Aggregation, Correlation, or Inference (P5) recognizes that individually non-sensitive
datasets may collectively reveal sensitive information when combined or interpreted 
together. Persistence Beyond Intended Lifetime (P6) addresses privacy risks
associated with temporary files, metadata snapshots, cached datasets, 
archived copies, and other artifacts that remain accessible after their 
intended purpose has ended.

The remaining threat category Misplaced Trust in Privacy by Format (P7) reflects broader
characteristics of modern HDF5 ecosystems. Examples include: 

\begin{itemize}

\item Misinterpretation of data privacy in HDF5 format leads to the possibility of 
sensitive information disclosure through metadata, structural information, or inference.

\item Conversion to HDF5 highlights that transforming legacy scientific formats 
into HDF5 often increases accessibility and interoperability by making implicit 
information explicit, but this additional semantic information may itself introduce new privacy risks. 

\item Loss of Data Control describes situations in which data movement, replication, 
backup, cloud migration, or transfer across organizational boundaries weakens the 
original privacy protections and governance mechanisms applied to the data.

\end{itemize}

Although presented as distinct categories, these mechanisms can occur together.
A single scientific workflow may enrich metadata during format conversion, aggregate 
datasets from multiple sources, generate temporary processing artifacts, and 
replicate results across multiple storage systems. Consequently, privacy risks should 
be evaluated across the complete scientific data lifecycle rather than 
by considering individual HDF5 files in isolation.

\paragraph{Recommendations}\mbox{}\\

Because HDF5 was designed to maximize scientific data sharing, portability, 
interoperability, and long-term usability rather than enforce privacy, 
no single mechanism can eliminate all privacy risks. Effective privacy protection 
therefore requires a multi-facet approach that considers the entire scientific data 
lifecycle, including data creation, processing, storage, sharing, transfer, 
and archival. Privacy should be treated as a property of the complete HDF5 ecosystem 
including applications, plugins, workflows, storage systems, and organizational 
practices rather than solely of the HDF5 file format. 
The following \code{CORE-PRIV-S} identifiers identify recommendations in this
track. Their numbers do not express risk, urgency, effort, dependency, or
implementation order.

For a confirmed \code{CORE-PRIV-001} sharing boundary, the Near-term route
applies: \code{RM-1D} first bounds the population and recipients and tests
minimization and recipient controls; \code{RM-2G} is conditional when policy
still requires whole-file or equivalent all-metadata protection.
\code{CORE-PRIV-S2} and extension controls support the unrated
\code{EXT-PRIV-001} evidence and assurance track and inherit no urgency from
\code{CORE-PRIV-001}.

\paragraph{Recommendation CORE-PRIV-S1: \emph{Privacy exposure awareness}}\mbox{}\\*
Scientists, software developers, data stewards, and infrastructure providers should
recognize that HDF5's self-describing nature inherently exposes contextual
information in many ways. Privacy reviews should therefore evaluate not only
scientific datasets, but also metadata, structural organization, workflow artifacts,
derived products, and information that may be inferred through aggregation or
correlation. Reviews should also take into consideration possible privacy
exposure during data lifecycle.

Awareness is not itself a control, so this recommendation is written to
produce an artifact rather than a disposition of mind.

\noindent\textbf{Acceptance evidence.}
A dated privacy review record exists for the defined data population and
sharing boundary, and it contains, at minimum:

\begin{enumerate}[label=(\alph*),leftmargin=2em]
    \item the named data owner and the intended recipient set for each
          disclosure boundary;
    \item an enumeration of sensitive items found in each asset class
          identified in Section~\ref{sec:core-lib-format} --- data,
          user-defined metadata, structural metadata, derived or reconstructed
          information, and workflow artifacts --- or a positive statement that
          the class was examined and found empty;
    \item the lifecycle stages examined, with the stages not examined named
          rather than omitted; and
    \item a disposition for every enumerated item: retain, minimize,
          tokenize, restrict recipients, or accept with a recorded rationale.
\end{enumerate}

A review that examines datasets only, or that records no negative findings for
the metadata and artifact classes, does not satisfy this recommendation.

\paragraph{Recommendation CORE-PRIV-S2: \emph{Privacy exposure through HDF5 tools and plugins awareness}}\mbox{}\\*
The HDF5 ecosystem extends beyond the core HDF5 software (i.e., file format and the library), 
and includes a diverse collection of command-line utilities, visualization tools, 
analysis libraries, VOL connectors, VFDs, filters, and third-party applications. 
These components can generate additional outputs such as logs, diagnostic
messages, temporary files, metadata snapshots, converted datasets, and cached 
intermediate results that may expose sensitive information beyond the contents of the 
original HDF5 file. Furthermore, plugins may introduce new metadata, or transform 
data in ways that alter their privacy characteristics. Consequently, privacy 
assessments should evaluate not only HDF5 files but also the behavior of the
software ecosystem that processes them. Developers should document the privacy
implications of their tools, minimize unnecessary generation of sensitive artifacts,
and provide users with mechanisms to control logging, metadata generation, and
retention of intermediate products.

\noindent\textbf{Acceptance evidence.}
For each in-scope tool, connector, filter, and VFD:

\begin{enumerate}[label=(\alph*),leftmargin=2em]
    \item published documentation lists the outputs it can generate --- stdout
          and stderr, redirected files, logs, temporary files, caches,
          snapshots, converted or derived files --- and the option, property,
          or environment control that suppresses or bounds each one;
    \item the list is confirmed by a reproducible trace of a representative
          invocation rather than asserted from the source, and the trace is
          retained; and
    \item any output observed in the trace but absent from the documentation
          is recorded as a defect against that component.
\end{enumerate}

Table~\ref{tab:hdf5-tools-privacy} is the starting inventory for the official
command-line tools, not a completed trace. This recommendation supports the
unrated \code{EXT-PRIV-001} evidence track and inherits no response urgency
from \code{CORE-PRIV-001}.

\paragraph{Recommendation CORE-PRIV-S3: \emph{Privacy-aware system design}}\mbox{}\\*
HDF5 software systems should be built with data privacy in mind. Applications should 
minimize collection and retention of unnecessary metadata, avoid embedding sensitive 
information in object names or descriptive attributes, remove temporary processing
artifacts when they are no longer needed, and carefully evaluate automatically
generated metadata introduced during data conversion or interoperability.
Workflow designers should review intermediate products, logs, cached datasets, and
derived files before publication or long-term archival.

\noindent\textbf{Acceptance evidence.}
The claim to be tested is that minimization actually removed the information,
not that it was removed from the API view:

\begin{enumerate}[label=(\alph*),leftmargin=2em]
    \item before-and-after inspection of a representative file through
          \emph{both} the HDF5 API and a raw-byte search --- at minimum
          \code{strings}-class extraction over the whole file, including the
          user block and any free or unreferenced space --- shows that each
          item marked \emph{minimize} or \emph{tokenize} in the
          \code{CORE-PRIV-S1} review is absent from both views;
    \item metadata introduced automatically during conversion is enumerated
          for the tested pipeline and each item is dispositioned; and
    \item temporary artifacts are removed under a stated retention rule, and
          the rule is verified by inspecting the working area after both
          normal and abnormal termination.
\end{enumerate}

An item that disappears from \code{h5dump} output but is still recoverable
from the raw bytes has not been minimized. Where the file is rewritten to
achieve this, record pre- and post-checksums and preserve provenance.

\paragraph{Recommendation CORE-PRIV-S4: \emph{Data encryption as privacy control}}\mbox{}\\*
Encryption represents an important but limited privacy control. 

Currently, HDF5 doesn’t support data encryption that allows I/O on encrypted 
data, effectively disabling sensitive data exposure during data 
transfer and storage. Users may implement custom raw data encryption 
and partial metadata encryption (i.e., names of the objects) using 
existing HDF5 filtering mechanism. Such approach still leaves user-defined 
metadata (e.g., attributes) and structural metadata unencrypted. Information 
that remains unencrypted as application-specific workflow artifacts, 
temporary files, system logs, or derived datasets may continue to 
disclose sensitive information.
  
HDF5 encryption feature that properly implements encryption of all data 
stored in HDF5 while allowing partial I/O on encrypted data is currently 
under development. However, encryption does not eliminate all HDF5 privacy risks. 
Once encrypted data are decrypted for legitimate processing, applications, 
plugins, caches, diagnostic logs, and intermediate workflow products may again 
become sources of unintended exposure. Encryption therefore reduces the attack 
surface but does not address inference, aggregation, excessive metadata 
generation, data retention, or loss of governance. Encryption also complicates 
workflows and may lead to data loss if on some stage of data life cycle 
encryption keys become unavailable.

Therefore, encryption should be viewed as one component of a comprehensive 
privacy strategy rather than as a complete solution. Effective protection 
requires combining encryption with careful metadata management, lifecycle-aware
data governance, privacy-aware application design, access controls provided by
the underlying computing environment, and education of users regarding the
unique privacy characteristics of HDF5-based scientific data systems.

\noindent\textbf{Acceptance evidence.}
Encryption is credited only for the boundaries actually tested, and only when
recoverability is demonstrated alongside confidentiality:

\begin{enumerate}[label=(\alph*),leftmargin=2em]
    \item ciphertext and metadata inspection shows that the items the
          \code{CORE-PRIV-S1} review marked as requiring protection are not
          observable at the protected boundary, tested through both the HDF5
          API and raw-byte search;
    \item an unauthorized-open test fails closed, and a partial-I/O workflow
          test confirms the protected file remains usable for its intended
          access pattern;
    \item a protected key inventory names a key custodian per key, and
          backup/restore evidence exists; and
    \item a loss-and-recovery exercise is performed and passes.
\end{enumerate}

Item (d) is not optional. Encryption without a tested recovery path converts a
disclosure risk into a durable data-loss risk, which is a different harm and
not an improvement. Until all four pass, the affected sharing boundary is
avoided or a tested access restriction is maintained within the
\code{CORE-PRIV-001} Near-term clock. This recommendation is the conditional
branch \code{RM-2G}, reached only where \code{RM-1D} minimization and
recipient controls cannot meet the disclosure policy.


\end{document}
